\section{USAMO 1979}
        \begin{enumerate}
        	\item (\textit{USAMO 1979}) Determine all non-negative integral solutions $(n_1,n_2,\dots , n_{14})$ if any, apart from permutations, of the Diophantine Equation $n_1^4+n_2^4+\cdots +n_{14}^4=1599$.


 

	\item (\textit{USAMO 1979}) $N$ is the north pole. $A$ and $B$ are points on a great circle through $N$ equidistant from $N$. $C$ is a point on the equator. Show that the great circle through $C$ and $N$ bisects the angle $ACB$ in the spherical triangle $ABC$ (a spherical triangle has great circle arcs as sides).


 

	\item (\textit{USAMO 1979}) $a_1, a_2, \ldots, a_n$ is an arbitrary sequence of positive integers. A member of the sequence is picked at 
random. Its value is $a$. Another member is picked at random, independently of the first. Its value is $b$. Then a third value, $c$. Show that the probability that $a + b + c$ is divisible by $3$ is at least $\frac14$.


 

	\item (\textit{USAMO 1979}) $P$ lies between the rays $OA$ and $OB$. Find $Q$ on $OA$ and $R$ on $OB$ collinear with $P$ so that $\frac{1}{PQ} + \frac{1}{PR}$ is as large as possible.


 

	\item (\textit{USAMO 1979}) Let $A_1,A_2,...,A_{n+1}$ be distinct subsets of $[n]$ with $|A_1|=|A_2|=\cdots =|A_{n+1}|=3$.  Prove that $|A_i\cap A_j|=1$ for some pair $\{i,j\}$. Note that $[n] = \{1, 2, 3, ..., n\}$, or, alternatively, $\{x: 1 \le x \le n\}$.


 

\end{enumerate}